\DocumentMetadata{
  pdfversion=2.0,
  pdfstandard=ua-2,
  lang=en-US,
  tagging=on,
  tagging-setup={math/setup=mathml-SE,tabsorder=structure}
}
\documentclass[11pt,letterpaper]{article}
\usepackage[margin=0.68in,includefoot]{geometry}
\usepackage{newcomputermodern}
\usepackage{amsmath,amscd,mathtools}
\providecommand{\square}{\mdlgwhtsquare}
\usepackage[hidelinks]{hyperref}
\hypersetup{
  pdftitle={Numerical Linear Algebra First-Year Exam, May 2025},
  pdfauthor={University of Florida Department of Mathematics},
  pdfsubject={Graduate mathematics examination},
  pdfdisplaydoctitle=true,
  bookmarksopen=true
}
\setcounter{secnumdepth}{0}
\setlength{\parindent}{0pt}
\setlength{\parskip}{0.28em}
\setlength{\emergencystretch}{3em}
\setlength{\leftmargini}{2.15em}
\setlength{\leftmarginii}{2.65em}
\setlength{\leftmarginiii}{3em}
\renewcommand{\labelenumi}{\textbf{\theenumi.}}
\pagestyle{empty}
\raggedbottom
\allowdisplaybreaks
\newenvironment{examproblems}[1]{%
  \begin{enumerate}%
  \setcounter{enumi}{#1}%
  \setlength{\topsep}{0.35em}%
  \setlength{\partopsep}{0pt}%
  \setlength{\itemsep}{0.48em}%
  \setlength{\parsep}{0.18em}%
}{\end{enumerate}}
\newenvironment{examsubparts}{%
  \begin{enumerate}%
  \setlength{\topsep}{0.18em}%
  \setlength{\partopsep}{0pt}%
  \setlength{\itemsep}{0.16em}%
  \setlength{\parsep}{0.1em}%
}{\end{enumerate}}
\newenvironment{examinstructions}{%
  \par\begingroup\itshape
}{%
  \par\endgroup\vspace{0.18em}
}
\makeatletter
\renewcommand\section{\@startsection{section}{1}{0pt}%
  {-1.25ex plus -.25ex minus -.1ex}%
  {0.55ex plus .1ex}%
  {\normalfont\large\bfseries}}
\renewcommand{\@maketitle}{%
  \newpage\null\vskip -1.2em%
  {\footnotesize\bfseries
    \href{https://www.ufl.edu/}{University of Florida}, \href{https://math.ufl.edu/}{Department of Mathematics}\par}%
  \vskip 0.18em%
  \tagstructbegin{tag=Title}%
    {\LARGE\bfseries\href{https://gma.math.ufl.edu/exams/numerical-linear-algebra-first-year/}{Numerical Linear Algebra First-Year Exam}, May 2025\par}%
  \tagstructend%
  \vskip 0.35em%
  \tagmcbegin{artifact}%
    \hrule height 1pt%
  \tagmcend%
  \vskip 0.55em%
}
\makeatother
\title{Numerical Linear Algebra First-Year Exam, May 2025}
\author{}
\date{}
\begin{document}
\maketitle
\thispagestyle{empty}
\begin{examinstructions}
Do 4 (four) problems.
\end{examinstructions}
\begin{examproblems}{0}
\item
Assume \(A \in \mathbb{C}^{m\times m}\).
\begin{examsubparts}
\item[\textnormal{(a)}]
Show that~\(A\) has a Schur decomposition.
\item[\textnormal{(b)}]
If~\(A\) is normal, show that~\(A\) is diagonalizable.
\end{examsubparts}
\item
Suppose~\(A\) is Hermitian positive definite.
\begin{examsubparts}
\item[\textnormal{(a)}]
Prove that each principal submatrix of~\(A\) is Hermitian positive definite.
\item[\textnormal{(b)}]
Prove that an element of~\(A\) with largest magnitude lies on the diagonal.
\item[\textnormal{(c)}]
Prove that~\(A\) has a Cholesky decomposition.
\end{examsubparts}
\item
This problem has two parts.
\begin{examsubparts}
\item[\textnormal{(a)}]
Show that \(\lVert x\rVert_\infty\) is equivalent to \(\lVert x\rVert_2\) for all \(x\in\mathbb{R}^n\). That is to find~\(C\) and~\(c\) such that \(c\lVert x\rVert_\infty\leq \lVert x\rVert_2\leq C\lVert x\rVert_\infty\), for all \(x\in\mathbb{R}^n\). Note that the constants should be determined so that the equalities hold for some nonzero \(x\in\mathbb{R}^n\).
\item[\textnormal{(b)}]
Show that \(\lVert QA\rVert_2=\lVert A\rVert_2\) if~\(Q\) is a unitary matrix.
\end{examsubparts}
\item
Assume that \(A\in\mathbb{C}^{n\times n}\) and there exists \(p\geq 1\) such that \(\lVert A\rVert_p<1\), where \(\lVert\cdot\rVert_p\) is a vector-induced matrix norm.
\begin{examsubparts}
\item[\textnormal{(a)}]
Prove that \(I-A\) is invertible.
\item[\textnormal{(b)}]
Prove that 
\[
(I-A)^{-1}=\sum_{k=0}^{\infty}A^k.
\]
\item[\textnormal{(c)}]
Prove that \(\lVert A\rVert_q\lVert A^{-1}\rVert_q\geq 1\), \(\forall\,1\leq q\leq\infty\).
\item[\textnormal{(d)}]
Prove that 
\[
\frac{1}{1+\lVert A\rVert_p}\leq \lVert(I-A)^{-1}\rVert_p\leq\frac{1}{1-\lVert A\rVert_p}.
\]
\end{examsubparts}
\item
Let \(A=U\Sigma V^*\) be the singular value decomposition of \(A\in\mathbb{C}^{m\times n}\). Let \(u_j\) denote column~\(j\) of~\(U\).
\begin{examsubparts}
\item[\textnormal{(a)}]
Suppose \(\operatorname{rank}(A)=p<n<m\). Show \(\{u_1,u_2,\ldots,u_p\}\) is a basis for \(\operatorname{Col}(A)\) and \(\{u_{p+1},u_{p+2},\ldots,u_m\}\) is a basis for \(\operatorname{Null}(A^*)\).
\item[\textnormal{(b)}]
Suppose~\(A\) is full rank and \(x\neq 0\). Let \(\sigma_i\), \(i=1,\ldots,n\), be the singular values of~\(A\). Show 
\[
\sigma_1\geq\frac{\lVert Ax\rVert_2}{\lVert x\rVert_2}\geq\sigma_n>0.
\]
 If you want to use the property that \(\lVert A\rVert_2=\sigma_1\), then you must prove that it holds.
\end{examsubparts}
\end{examproblems}
\end{document}
